\documentclass[11pt,letterpaper]{article}

% ------------------------------------------------------------
% Manual didáctico de uso: hidden-attractors-fo / version_2
% ------------------------------------------------------------

\usepackage[spanish,es-nodecimaldot]{babel}
\usepackage[utf8]{inputenc}
\usepackage[T1]{fontenc}
\usepackage{lmodern}
\usepackage{geometry}
\usepackage{microtype}
\usepackage{xcolor}
\usepackage{amsmath,amssymb,amsfonts}
\usepackage{booktabs}
\usepackage{array}
\usepackage{enumitem}
\usepackage{listings}
\usepackage{hyperref}
\usepackage{fancyhdr}
\usepackage{titlesec}
\usepackage{tcolorbox}
\usepackage{caption}

\geometry{margin=2.2cm}
\hypersetup{
  colorlinks=true,
  linkcolor=blue!45!black,
  urlcolor=blue!45!black,
  citecolor=blue!45!black,
  pdftitle={Manual didáctico de usuario de hidden-attractors-fo},
  pdfauthor={Proyecto Hidden Attractors Fractional Order}
}

\newcolumntype{L}[1]{>{\raggedright\arraybackslash}p{#1}}

\definecolor{codebg}{RGB}{248,248,248}
\definecolor{coderule}{RGB}{210,210,210}
\definecolor{warningbg}{RGB}{255,248,230}
\definecolor{infobg}{RGB}{235,244,255}
\definecolor{okbg}{RGB}{235,250,240}

\lstdefinestyle{terminal}{
  basicstyle=\ttfamily\small,
  backgroundcolor=\color{codebg},
  frame=single,
  rulecolor=\color{coderule},
  breaklines=true,
  columns=fullflexible,
  keepspaces=true,
  showstringspaces=false
}

\lstdefinestyle{python}{
  language=Python,
  basicstyle=\ttfamily\small,
  backgroundcolor=\color{codebg},
  frame=single,
  rulecolor=\color{coderule},
  breaklines=true,
  columns=fullflexible,
  keepspaces=true,
  showstringspaces=false,
  keywordstyle=\color{blue!55!black},
  commentstyle=\color{green!35!black},
  stringstyle=\color{red!45!black}
}

\newtcolorbox{advertencia}{
  colback=warningbg,
  colframe=orange!65!black,
  title=Advertencia técnica,
  fonttitle=\bfseries
}

\newtcolorbox{idea}{
  colback=infobg,
  colframe=blue!55!black,
  title=Idea operativa,
  fonttitle=\bfseries
}

\newtcolorbox{resultado}{
  colback=okbg,
  colframe=green!50!black,
  title=Qué debe aparecer si todo salió bien,
  fonttitle=\bfseries
}

\pagestyle{fancy}
\fancyhf{}
\lhead{Manual de usuario: \texttt{hidden-attractors-fo}}
\rhead{\thepage}
\setlength{\headheight}{14pt}

\titleformat{\section}{\Large\bfseries}{\thesection.}{0.6em}{}
\titleformat{\subsection}{\large\bfseries}{\thesubsection.}{0.6em}{}
\titleformat{\subsubsection}{\normalsize\bfseries}{\thesubsubsection.}{0.6em}{}

\title{\textbf{Manual didáctico de usuario de la librería}\\[2mm]
\texttt{hidden-attractors-fo}\\[1mm]
\large Guía explícita para ejecutar, interpretar y defender los workflows de \texttt{version\_2}}
\author{Proyecto: localización de atractores ocultos en sistemas fraccionarios}
\date{Versión de trabajo: \today}

\begin{document}
\maketitle
\tableofcontents
\newpage

% ============================================================
\section{Propósito de este manual}

Este documento explica cómo usar la carpeta \texttt{version\_2} del repositorio como librería Python. No es una descripción del backend en abstracto. El objetivo es responder, de forma operativa, las siguientes preguntas:

\begin{enumerate}[label=\arabic*]
  \item ¿Dónde debo pararme en la terminal?
  \item ¿Qué debo instalar?
  \item ¿Qué comando corro primero?
  \item ¿Qué calcula cada comando?
  \item ¿Qué archivos de salida debo revisar?
  \item ¿Cómo sé si el resultado apoya o descarta la ocultedad?
  \item ¿Qué partes son librería pública y qué partes siguen siendo scripts heredados?
\end{enumerate}

\begin{advertencia}
La librería no demuestra por sí sola la existencia matemática de atractores ocultos. Produce evidencia numérica bajo contratos explícitos: modelo, orden fraccionario $q$, paso $h$, memoria finita $L_m$, tiempo final, tiempo de descarte, backend y umbrales de clasificación. Una conclusión de ocultedad requiere verificar que la cuenca del atractor candidato no intersecta vecindades de ningún equilibrio bajo las pruebas realizadas.
\end{advertencia}

% ============================================================
\section{Qué es exactamente \texttt{version\_2}}

La carpeta \texttt{version\_2} es la versión organizada como paquete Python instalable. Su paquete se llama:

\begin{center}
\texttt{hidden-attractors-fo}
\end{center}

y el módulo importable se llama:

\begin{center}
\texttt{hidden\_attractors}
\end{center}

La idea central es que el código ya no debe usarse como una colección desordenada de scripts sueltos, sino como una librería con módulos reutilizables:

\begin{itemize}
  \item modelos dinámicos;
  \item carga de candidatos finales;
  \item diagnósticos de trayectorias;
  \item etiquetas de cuencas;
  \item gráficas;
  \item workflows de robustez y verificación;
  \item envoltorios hacia backends nativos en C/EFORK.
\end{itemize}

\subsection{Límite entre librería y scripts heredados}

La regla práctica es:

\begin{center}
\begin{tabular}{p{0.38\linewidth}p{0.52\linewidth}}
\toprule
Carpeta & Cómo debes usarla \\
\midrule
\texttt{hidden\_attractors/} & API principal de la librería. Aquí está lo que se debe importar desde Python. \\
\texttt{examples/} & Ejemplos cortos. Sirven para aprender y verificar instalación. \\
\texttt{tools/cli/} & Comandos mantenidos para ejecutar workflows desde terminal. \\
\texttt{docs/} & Documentación técnica y reportes. \\
\texttt{outputs/} & Salidas de referencia y nuevas salidas generadas por corridas. \\
\texttt{tools/legacy/} & Scripts históricos. Pueden contener lógica útil, pero no deben crecer como API nueva. \\
\bottomrule
\end{tabular}
\end{center}

\begin{idea}
Para uso diario, empieza con \texttt{examples/} y \texttt{tools/cli/}. Solo después entra a \texttt{hidden\_attractors/} si necesitas programar tus propios análisis.
\end{idea}

% ============================================================
\section{Instalación desde cero}

\subsection{Ubicación correcta en la terminal}

Primero entra al repositorio y luego a la carpeta \texttt{version\_2}. En Windows PowerShell:

\begin{lstlisting}[style=terminal]
cd C:\Users\moren\Desktop\Hidden-Attractors-Localization\version_2
\end{lstlisting}

En Linux/macOS:

\begin{lstlisting}[style=terminal]
cd ~/Hidden-Attractors-Localization/version_2
\end{lstlisting}

El archivo \texttt{pyproject.toml} está dentro de \texttt{version\_2}. Por eso, si ejecutas \texttt{pip install -e .} desde la raíz equivocada, Python puede no encontrar el paquete.

\subsection{Crear entorno virtual recomendado}

En PowerShell:

\begin{lstlisting}[style=terminal]
python -m venv .venv
.\.venv\Scripts\Activate.ps1
python -m pip install --upgrade pip
\end{lstlisting}

En Linux/macOS:

\begin{lstlisting}[style=terminal]
python -m venv .venv
source .venv/bin/activate
python -m pip install --upgrade pip
\end{lstlisting}

\subsection{Instalar la librería en modo editable}

Instalación mínima:

\begin{lstlisting}[style=terminal]
python -m pip install -e .
\end{lstlisting}

Instalación para desarrollo y pruebas:

\begin{lstlisting}[style=terminal]
python -m pip install -e ".[dev]"
\end{lstlisting}

Instalación con análisis externos opcionales:

\begin{lstlisting}[style=terminal]
python -m pip install -e ".[analysis]"
\end{lstlisting}

Instalación para construir documentación con MkDocs:

\begin{lstlisting}[style=terminal]
python -m pip install -e ".[docs]"
\end{lstlisting}

\subsection{Dependencias mínimas}

La instalación base usa principalmente:

\begin{itemize}
  \item Python $\geq 3.10$;
  \item \texttt{numpy};
  \item \texttt{matplotlib};
  \item compilador C si se ejecutan workflows nativos;
  \item \texttt{pytest} solo si instalas \texttt{.[dev]}.
\end{itemize}

En Windows, para workflows nativos, necesitas que \texttt{gcc} esté en el \texttt{PATH}. En macOS puede ser necesario instalar soporte OpenMP mediante \texttt{libomp}.

% ============================================================
\section{Prueba mínima de instalación}

Después de instalar, ejecuta:

\begin{lstlisting}[style=terminal]
python -m compileall hidden_attractors examples tests tools/cli
\end{lstlisting}

Luego:

\begin{lstlisting}[style=terminal]
python examples/quickstart_equilibria.py
python examples/list_final_candidates.py
python examples/dynamical_analysis_gallery.py
\end{lstlisting}

Si instalaste con \texttt{.[dev]}, ejecuta:

\begin{lstlisting}[style=terminal]
python -m pytest -q
\end{lstlisting}

\begin{resultado}
La prueba de equilibrios debe imprimir tres puntos: \texttt{E0}, \texttt{E+} y \texttt{E-}, con norma residual cercana a cero. La prueba de candidatos debe listar tres candidatos finales. El comando de \texttt{pytest} debe pasar las pruebas de humo.
\end{resultado}

% ============================================================
\section{Modelo mental de la librería}

La librería está organizada alrededor del siguiente flujo:

\begin{center}
\begin{tabular}{p{0.22\linewidth}p{0.68\linewidth}}
\toprule
Etapa & Pregunta que responde \\
\midrule
Modelo & ¿Qué sistema dinámico se está usando? \\
Equilibrios & ¿Cuáles son los puntos de equilibrio que deben excluirse de la cuenca del atractor candidato? \\
Candidatos & ¿Qué semillas o puntos finales de continuación se van a verificar? \\
Trayectorias & ¿Qué hace el sistema de Caputo desde una condición inicial? \\
Robustez & ¿El comportamiento observado persiste al cambiar $h$, $L_m$ o el tiempo final? \\
Controles esféricos & ¿La cuenca candidata toca vecindades pequeñas de algún equilibrio? \\
Cuencas refinadas & ¿Qué pasa con puntos previamente marcados como \texttt{unknown}? \\
Reporte & ¿Qué archivos justifican numéricamente la conclusión? \\
\bottomrule
\end{tabular}
\end{center}

\subsection{Qué no hace automáticamente}

\begin{advertencia}
\texttt{version\_2} no debe interpretarse como una máquina que, al ejecutarse, encuentra y demuestra atractores ocultos. La generación histórica de semillas por Lur'e/Nyquist, función descriptiva clásica o Machado/FDF aparece en scripts heredados o salidas de referencia. La librería activa se concentra en cargar candidatos, analizarlos, simularlos, clasificarlos y generar evidencia reproducible.
\end{advertencia}

% ============================================================
\section{El sistema de Chua implementado}

El modelo base en \texttt{hidden\_attractors.models.chua} usa la no linealidad por tramos:

\begin{equation}
 f(x)=m_1x+\frac{1}{2}(m_0-m_1)\bigl(|x+1|-|x-1|\bigr).
\end{equation}

El campo vectorial se escribe como:

\begin{align}
{}^C D_t^q x &= \alpha\bigl(y-x-f(x)\bigr),\\
{}^C D_t^q y &= x-y+z,\\
{}^C D_t^q z &= -\beta y-\gamma z.
\end{align}

En la librería, la función \texttt{rhs\_piecewise} no integra el sistema. Solo evalúa el lado derecho:

\begin{equation}
F(X)=
\begin{bmatrix}
\alpha(y-x-f(x))\\
x-y+z\\
-\beta y-\gamma z
\end{bmatrix}.
\end{equation}

La integración fraccionaria causal de Caputo se delega al backend nativo EFORK cuando se ejecutan workflows que simulan trayectorias largas.

\subsection{Parámetros por defecto}

La clase \texttt{ChuaParameters} usa por defecto:

\begin{center}
\begin{tabular}{cc}
\toprule
Parámetro & Valor \\
\midrule
$\alpha$ & $8.4562$ \\
$\beta$ & $12.0732$ \\
$\gamma$ & $0.0052$ \\
$m_0$ & $-0.1768$ \\
$m_1$ & $-1.1468$ \\
\bottomrule
\end{tabular}
\end{center}

\subsection{Equilibrios}

La función:

\begin{lstlisting}[style=python]
from hidden_attractors.models import equilibria_piecewise
\end{lstlisting}

calcula los tres equilibrios del sistema por tramos:

\begin{equation}
E_0=(0,0,0), \qquad E_+, \qquad E_-.
\end{equation}

Para los equilibrios exteriores se usa la relación algebraica implementada en la librería. Si

\begin{equation}
\kappa=-\frac{\beta}{\gamma+\beta},
\end{equation}

entonces

\begin{equation}
 x_+ = -\frac{m_0-m_1}{m_1-\kappa},
 \qquad
 x_- = \frac{m_0-m_1}{m_1-\kappa}.
\end{equation}

Después, para cada $x$, la librería reconstruye:

\begin{equation}
E(x)=\bigl(x,\,x+f(x),\,f(x)\bigr).
\end{equation}

\begin{advertencia}
Calcular los equilibrios no clasifica estabilidad fraccionaria. Para estabilidad local del sistema de orden $q$ debe revisarse el espectro del Jacobiano y aplicar el criterio de Matignon:
\begin{equation}
|\arg(\lambda_i)|>\frac{q\pi}{2}.
\end{equation}
La librería, en esta parte, solo da las coordenadas de equilibrio y verifica residuales del campo vectorial.
\end{advertencia}

% ============================================================
\section{Primer uso: calcular y verificar equilibrios}

Ejecuta:

\begin{lstlisting}[style=terminal]
python examples/quickstart_equilibria.py
\end{lstlisting}

Este script hace internamente:

\begin{lstlisting}[style=python]
from hidden_attractors import chua_piecewise_parameters
from hidden_attractors.models import equilibria_piecewise, rhs_piecewise

params = chua_piecewise_parameters()
for name, point in equilibria_piecewise(params).items():
    residual_norm = norm(rhs_piecewise(point, params))
    print(name, point, residual_norm)
\end{lstlisting}

\subsection{Qué calcula}

Calcula las coordenadas de $E_0$, $E_+$ y $E_-$ y evalúa:

\begin{equation}
\|F(E_i)\|.
\end{equation}

Si el residual es casi cero, el punto es efectivamente equilibrio del campo vectorial implementado.

\subsection{Cómo interpretar la salida}

Una salida esperada tiene la forma:

\begin{lstlisting}[style=terminal]
E0: point=[0.0, 0.0, 0.0] residual_norm=0.000e+00
E+: point=[...] residual_norm=...e-16
E-: point=[...] residual_norm=...e-16
\end{lstlisting}

Si aparece un residual grande, hay una inconsistencia de parámetros, de implementación o de entorno.

% ============================================================
\section{Segundo uso: listar candidatos finales}

Ejecuta:

\begin{lstlisting}[style=terminal]
python examples/list_final_candidates.py
\end{lstlisting}

O, después de instalar la librería:

\begin{lstlisting}[style=terminal]
hidden-attractors-list-candidates
\end{lstlisting}

\subsection{Qué calcula}

No simula. No integra. No verifica ocultedad. Solo carga registros de candidatos desde salidas de referencia guardadas en \texttt{outputs/}. Cada candidato se representa con la clase \texttt{CandidateRecord}.

Cada registro contiene, cuando está disponible:

\begin{center}
\begin{tabular}{p{0.28\linewidth}p{0.62\linewidth}}
\toprule
Campo & Significado \\
\midrule
\texttt{candidate\_id} & Identificador único del candidato. \\
\texttt{route} & Ruta por la cual se obtuvo: por ejemplo, \texttt{Machado\_FDF} o \texttt{Lure\_rank\_0001}. \\
\texttt{q} & Orden fraccionario usado en la corrida de referencia. \\
\texttt{seed} & Semilla inicial teórica o numérica asociada al candidato. \\
\texttt{robust\_start} & Punto desde donde se inician pruebas de robustez. Suele venir del final de una continuación o verificación previa. \\
\texttt{mu}, \texttt{theta} & Parámetros de la ruta Machado/FDF, si aplica. \\
\texttt{A}, \texttt{sigma0}, \texttt{omega} & Datos asociados al balance armónico o ruta Lur'e, si existen. \\
\texttt{rho\_H}, \texttt{residual\_abs} & Indicadores numéricos de la ruta de búsqueda, si existen. \\
\texttt{source} & Archivo de salida desde donde se cargó el registro. \\
\bottomrule
\end{tabular}
\end{center}

\subsection{Cómo interpretar la salida}

Una línea típica tiene la forma:

\begin{lstlisting}[style=terminal]
branch_0_mu_4p00000_theta_0p00000 | route=Machado_FDF | q=0.9998 | start=[...] | seed=[...]
\end{lstlisting}

Esto significa que la librería encontró un registro de candidato ya existente. No significa que el candidato sea oculto.

\begin{advertencia}
Un \texttt{CandidateRecord} es una hipótesis inicial. La ocultedad se evalúa después, con pruebas de cuenca alrededor de todos los equilibrios.
\end{advertencia}

% ============================================================
\section{Tercer uso: cargar una trayectoria y calcular métricas}

La librería usa la convención de trayectoria:

\begin{equation}
\texttt{t,x,y,z}.
\end{equation}

Es decir, un CSV de trayectoria debe tener columnas:

\begin{lstlisting}[style=terminal]
t,x,y,z
0.00,...,...,...
0.01,...,...,...
...
\end{lstlisting}

\subsection{Cargar trayectoria}

\begin{lstlisting}[style=python]
from hidden_attractors.io import load_trajectory_csv

traj = load_trajectory_csv("outputs/mi_corrida/trajectories/caso.csv")
\end{lstlisting}

\subsection{Calcular métricas}

\begin{lstlisting}[style=python]
from hidden_attractors.analysis import trajectory_metrics

metrics, payload = trajectory_metrics(
    traj,
    h=0.01,
    t_start=750.0,
    divergence_norm=120.0,
    equilibrium_tol=1e-3,
)

print(metrics)
\end{lstlisting}

\subsection{Qué calcula \texttt{trajectory\_metrics}}

La función calcula diagnósticos de tiempo finito:

\begin{center}
\begin{tabular}{L{0.28\linewidth}L{0.62\linewidth}}
\toprule
Métrica & Interpretación \\
\midrule
\texttt{bounded} & La trayectoria permanece finita y no supera la norma de divergencia definida. \\
\texttt{diverged} & La trayectoria diverge o contiene valores no finitos. \\
\texttt{equilibrium\_like} & El punto final cae cerca de un equilibrio, según \texttt{equilibrium\_tol}. \\
\texttt{noncollapsed\_variance} & La cola de la trayectoria conserva varianza suficiente; evita confundir un colapso numérico con oscilación. \\
\texttt{final\_norm} & Norma del estado final. \\
\texttt{max\_norm} & Norma máxima observada. \\
\texttt{range\_x}, \texttt{range\_y}, \texttt{range\_z} & Amplitud observada en la cola temporal. \\
\texttt{var\_x\_tail}, \texttt{var\_y\_tail}, \texttt{var\_z\_tail} & Varianza en la cola. \\
\texttt{fft\_peak} & Frecuencia dominante estimada por FFT para $x(t)$. \\
\texttt{psd\_entropy} & Entropía espectral normalizada aproximada. \\
\texttt{section\_points} & Número de cruces ascendentes de la sección $x=0$. \\
\texttt{cloud\_median\_distance} & Distancia mediana entre nubes de puntos, si se compara contra una referencia. \\
\texttt{section\_median\_distance} & Distancia entre nubes de sección de Poincaré, si existe referencia. \\
\bottomrule
\end{tabular}
\end{center}

\begin{advertencia}
Estas métricas no sustituyen exponentes de Lyapunov ni demuestran caos. Sirven para comparar trayectorias, detectar divergencia, detectar convergencia a equilibrio y medir similitud geométrica con una referencia.
\end{advertencia}

% ============================================================
\section{Cuarto uso: generar gráficas didácticas}

Para probar el módulo de análisis visual:

\begin{lstlisting}[style=terminal]
python examples/dynamical_analysis_gallery.py
\end{lstlisting}

Con una trayectoria específica:

\begin{lstlisting}[style=terminal]
python examples/dynamical_analysis_gallery.py --trajectory-csv outputs/mi_corrida/trajectories/caso.csv
\end{lstlisting}

\subsection{Qué produce}

El script puede generar:

\begin{itemize}
  \item retrato de fase 3D;
  \item proyecciones $xy$, $xz$, $yz$;
  \item series temporales;
  \item puntos de bifurcación extraídos de trayectorias existentes;
  \item \texttt{summary.json};
  \item \texttt{bifurcation\_points.csv}.
\end{itemize}

\subsection{Qué no produce}

No hace continuación real. No encuentra ramas de bifurcación. Solo postprocesa trayectorias ya generadas.

% ============================================================
\section{Flujo completo recomendado}

El flujo recomendado para usar la librería sin perderse es:

\begin{center}
\begin{tabular}{L{0.08\linewidth}L{0.32\linewidth}L{0.48\linewidth}}
\toprule
Paso & Acción & Comando principal \\
\midrule
1 & Instalar librería & \texttt{python -m pip install -e .} \\
2 & Probar importación & \texttt{python -m compileall hidden\_attractors examples tests tools/cli} \\
3 & Verificar equilibrios & \texttt{python examples/quickstart\_equilibria.py} \\
4 & Listar candidatos & \texttt{python examples/list\_final\_candidates.py} \\
5 & Revisar ayuda de workflows & \texttt{hidden-attractors-robustness-overlay --help} \\
6 & Ejecutar robustez geométrica & \texttt{hidden-attractors-robustness-overlay} \\
7 & Ejecutar controles alrededor de equilibrios & \texttt{hidden-attractors-sphere-controls} \\
8 & Refinar cuencas desconocidas, si existe una cuenca gruesa previa & \texttt{hidden-attractors-refined-basin} \\
9 & Revisar CSV/JSON finales & Abrir \texttt{summary.json}, \texttt{decision.csv} y gráficas. \\
\bottomrule
\end{tabular}
\end{center}

\begin{idea}
No empieces corriendo el workflow más pesado. Primero verifica equilibrios y candidatos. Después ejecuta \texttt{--help}. Solo entonces lanza simulaciones largas.
\end{idea}

% ============================================================
\section{Workflow 1: \texttt{robustness\_overlay}}

\subsection{Comando}

Después de instalar:

\begin{lstlisting}[style=terminal]
hidden-attractors-robustness-overlay
\end{lstlisting}

O directamente desde el wrapper:

\begin{lstlisting}[style=terminal]
python tools/cli/robustness_overlay_c_trajectories.py
\end{lstlisting}

Para ver opciones sin ejecutar simulaciones:

\begin{lstlisting}[style=terminal]
hidden-attractors-robustness-overlay --help
\end{lstlisting}

\subsection{Propósito}

Este workflow compara trayectorias iniciadas desde los candidatos finales al cambiar:

\begin{itemize}
  \item paso de integración $h$;
  \item longitud de memoria finita $L_m$;
  \item tiempo final de integración;
  \item tiempo de descarte transitorio.
\end{itemize}

Su pregunta es:

\begin{center}
\textit{¿El atractor candidato conserva geometría y rasgos espectrales bajo perturbaciones numéricas razonables?}
\end{center}

\subsection{Qué no responde}

No responde si el atractor es oculto. Para ocultedad se necesitan pruebas alrededor de los equilibrios.

\subsection{Casos de robustez por defecto}

La librería define casos similares a:

\begin{center}
\begin{tabular}{ccccc}
\toprule
Caso & $q$ & $h$ & $L_m$ & $T$ \\
\midrule
\texttt{R0\_base} & $0.9998$ & $0.01$ & $10$ & $1500$ \\
\texttt{R1\_h\_finer} & $0.9998$ & $0.005$ & $10$ & $1500$ \\
\texttt{R2\_h\_coarser} & $0.9998$ & $0.02$ & $10$ & $1500$ \\
\texttt{R3\_Lm\_lower} & $0.9998$ & $0.01$ & $5$ & $1500$ \\
\texttt{R4\_Lm\_higher} & $0.9998$ & $0.01$ & $20$ & $1500$ \\
\texttt{R5\_t\_longer} & $0.9998$ & $0.01$ & $10$ & $3000$ \\
\bottomrule
\end{tabular}
\end{center}

\subsection{Ejecución con carpeta de salida explícita}

Recomendado:

\begin{lstlisting}[style=terminal]
hidden-attractors-robustness-overlay --output-dir outputs/manual_run/robustness_overlay_test
\end{lstlisting}

Así no pierdes la corrida en una carpeta con timestamp.

\subsection{Archivos que genera}

\begin{center}
\begin{tabular}{L{0.38\linewidth}L{0.52\linewidth}}
\toprule
Archivo o carpeta & Significado \\
\midrule
\texttt{robustness\_overlay\_config.json} & Contrato numérico completo: candidatos, casos, parámetros y umbrales. \\
\texttt{launch\_manifest.json} & Procesos lanzados, PID y comandos usados. \\
\texttt{logs/} & Salidas \texttt{.out} y errores \texttt{.err} de cada proceso. \\
\texttt{trajectories/<candidate>/R*.csv} & Trayectorias muestreadas para cada candidato y caso de robustez. \\
\texttt{metrics\_<candidate>.csv} & Métricas individuales por candidato. \\
\texttt{metrics\_<candidate>.done} & Marca de que ese candidato terminó. \\
\texttt{robustness\_overlay\_metrics.csv} & Tabla agregada de métricas de todos los candidatos. \\
\texttt{robustness\_overlay\_summary.json} & Resumen global de la corrida. \\
\texttt{plots/overlay\_<candidate>.png} & Comparación visual de trayectorias. \\
\bottomrule
\end{tabular}
\end{center}

\subsection{Cómo saber si terminó bien}

Abre:

\begin{lstlisting}[style=terminal]
outputs/manual_run/robustness_overlay_test/robustness_overlay_summary.json
\end{lstlisting}

Debe contener:

\begin{lstlisting}[style=terminal]
"status": "ok"
\end{lstlisting}

Si dice \texttt{partial}, revisa:

\begin{lstlisting}[style=terminal]
outputs/manual_run/robustness_overlay_test/logs/
\end{lstlisting}

y busca archivos \texttt{.err} no vacíos.

\subsection{Cómo interpretar \texttt{robustness\_overlay\_metrics.csv}}

Columnas importantes:

\begin{center}
\begin{tabular}{L{0.34\linewidth}L{0.56\linewidth}}
\toprule
Columna & Lectura práctica \\
\midrule
\texttt{candidate\_id} & Candidato evaluado. \\
\texttt{route} & Ruta de origen del candidato. \\
\texttt{case\_id} & Caso de robustez. \\
\texttt{bounded} & Debe ser \texttt{True} para conservar interés. \\
\texttt{diverged} & Si es \texttt{True}, ese caso falló por divergencia. \\
\texttt{equilibrium\_like} & Si es \texttt{True}, la trayectoria terminó cerca de un equilibrio. \\
\texttt{noncollapsed\_variance} & Si es \texttt{False}, la trayectoria perdió dinámica relevante. \\
\texttt{range\_relative\_distance} & Distancia relativa de rangos respecto al caso base. Menor suele indicar mayor robustez geométrica. \\
\texttt{fft\_relative\_delta} & Cambio relativo de frecuencia dominante respecto al caso base. \\
\texttt{cloud\_median\_distance\_norm} & Distancia normalizada entre nubes de trayectoria. \\
\texttt{section\_median\_distance\_norm} & Distancia normalizada entre secciones de Poincaré. \\
\bottomrule
\end{tabular}
\end{center}

\begin{advertencia}
Un candidato robusto bajo \texttt{robustness\_overlay} todavía puede ser autoexcitado si alguna trayectoria iniciada cerca de un equilibrio llega al mismo atractor.
\end{advertencia}

% ============================================================
\section{Workflow 2: \texttt{sphere\_controls}}

\subsection{Comando}

\begin{lstlisting}[style=terminal]
hidden-attractors-sphere-controls
\end{lstlisting}

O:

\begin{lstlisting}[style=terminal]
python tools/cli/lure_top3_sphere_robustness.py
\end{lstlisting}

Para revisar opciones:

\begin{lstlisting}[style=terminal]
hidden-attractors-sphere-controls --help
\end{lstlisting}

\subsection{Propósito}

Este workflow sí está relacionado directamente con la verificación de ocultedad. Su pregunta es:

\begin{center}
\textit{¿La cuenca del atractor candidato intersecta vecindades pequeñas de algún equilibrio?}
\end{center}

Para responder, el workflow genera condiciones iniciales sobre esferas pequeñas alrededor de cada equilibrio:

\begin{equation}
X_0 = E_i + r\,u,
\end{equation}

donde $E_i$ es un equilibrio, $r$ es un radio pequeño y $u$ es una dirección unitaria.

Después clasifica cada trayectoria con el backend de cuencas.

\subsection{Interpretación matemática}

La definición operativa usada es:

\begin{quote}
Un atractor es oculto si su cuenca de atracción no intersecta ninguna vecindad abierta de los puntos de equilibrio.
\end{quote}

Por tanto:

\begin{itemize}
  \item si desde una esfera alrededor de un equilibrio aparecen trayectorias hacia el atractor candidato, hay evidencia contra la ocultedad;
  \item si no aparecen impactos hacia el atractor candidato bajo los radios y muestras probados, el resultado es compatible con ocultedad, pero no es una prueba absoluta.
\end{itemize}

\subsection{Ejecución recomendada}

Ejemplo con carpeta explícita:

\begin{lstlisting}[style=terminal]
hidden-attractors-sphere-controls --output-dir outputs/manual_run/sphere_controls_test
\end{lstlisting}

Ejemplo aumentando muestras y procesos:

\begin{lstlisting}[style=terminal]
hidden-attractors-sphere-controls ^
  --output-dir outputs/manual_run/sphere_controls_deep ^
  --samples-per-radius 1000 ^
  --chunks 8 ^
  --radii "1e-6,3e-6,1e-5,3e-5,1e-4,3e-4,1e-3"
\end{lstlisting}

En Linux/macOS cambia \texttt{\^{}} por \texttt{\\}:

\begin{lstlisting}[style=terminal]
hidden-attractors-sphere-controls \
  --output-dir outputs/manual_run/sphere_controls_deep \
  --samples-per-radius 1000 \
  --chunks 8 \
  --radii "1e-6,3e-6,1e-5,3e-5,1e-4,3e-4,1e-3"
\end{lstlisting}

\subsection{Parámetros importantes}

\begin{center}
\begin{tabular}{L{0.28\linewidth}L{0.62\linewidth}}
\toprule
Parámetro & Significado \\
\midrule
\texttt{--top-n} & Número de candidatos finales a verificar. Por defecto, 3. \\
\texttt{--radii} & Radios de las esferas alrededor de cada equilibrio. \\
\texttt{--samples-per-radius} & Número de direcciones por radio y equilibrio. \\
\texttt{--chunks} & Número de procesos independientes para dividir las muestras. \\
\texttt{--q} & Orden fraccionario. \\
\texttt{--h} & Paso de integración. \\
\texttt{--memory-length} & Longitud de memoria finita $L_m$. \\
\texttt{--t-final} & Tiempo final de integración. \\
\texttt{--t-burn} & Tiempo transitorio descartado. \\
\texttt{--divergence-norm} & Umbral de divergencia. \\
\texttt{--equilibrium-tol} & Tolerancia para considerar convergencia a equilibrio. \\
\texttt{--mean-x-gap} & Criterio operativo usado por el clasificador para separar clases objetivo. \\
\texttt{--cap-win} & Ventana usada por el clasificador para estimaciones de cola. \\
\bottomrule
\end{tabular}
\end{center}

\subsection{Archivos que genera}

\begin{center}
\begin{tabular}{L{0.38\linewidth}L{0.52\linewidth}}
\toprule
Archivo & Significado \\
\midrule
\path{top3_sphere_robustness_config.json} & Contrato completo de la prueba: candidatos, equilibrios, radios y parámetros. \\
\texttt{sphere\_plan.csv} & Todas las condiciones iniciales generadas sobre esferas. \\
\path{sphere_raw_chunk_*.csv} & Resultados por proceso. \\
\texttt{sphere\_raw.csv} & Resultados agregados de todas las muestras. \\
\path{sphere_cumulative_summary.csv} & Resumen acumulado por candidato, equilibrio y radio. \\
\texttt{sphere\_decision.csv} & Tabla de decisión principal sobre compatibilidad con ocultedad. \\
\path{attractor_robustness_raw.csv} & Clasificación de puntos finales de candidatos bajo contratos de robustez. \\
\path{attractor_robustness_summary.csv} & Resumen de robustez por candidato. \\
\path{top3_sphere_robustness_summary.json} & Resumen global de la corrida. \\
\texttt{logs/} & Salida y errores de cada proceso. \\
\bottomrule
\end{tabular}
\end{center}

\subsection{Cómo interpretar \texttt{sphere\_decision.csv}}

Columnas clave:

\begin{center}
\begin{tabular}{L{0.34\linewidth}L{0.56\linewidth}}
\toprule
Columna & Interpretación \\
\midrule
\texttt{candidate\_id} & Candidato evaluado. \\
\texttt{total\_sphere\_trajectories} & Total de trayectorias lanzadas desde vecindades de equilibrios. \\
\texttt{total\_target\_hits} & Número de trayectorias que llegaron a una clase objetivo. \\
\path{smallest_radius_with_target_hit} & Radio más pequeño donde se detectó contacto con la cuenca objetivo. \\
\texttt{hiddenness\_status} & Lectura operativa de la prueba. \\
\texttt{notes} & Recordatorio de que no es prueba absoluta. \\
\bottomrule
\end{tabular}
\end{center}

\subsection{Estados posibles de \texttt{hiddenness\_status}}

\begin{center}
\begin{tabular}{L{0.42\linewidth}L{0.48\linewidth}}
\toprule
Estado & Significado \\
\midrule
\path{not_supported_by_sphere_equilibrium_test} & Se encontró al menos una trayectoria desde vecindades de equilibrio hacia la clase objetivo. Esto contradice la hipótesis de ocultedad bajo el contrato probado. \\
\path{compatible_with_hiddenness_under_tested_spheres} & No se detectaron impactos a la clase objetivo desde las esferas probadas. Es compatible con ocultedad, pero no prueba definitiva. \\
\bottomrule
\end{tabular}
\end{center}

\begin{advertencia}
Si \texttt{total\_target\_hits > 0}, no conviene defender ese candidato como oculto con ese contrato numérico. Si \texttt{total\_target\_hits = 0}, se puede decir que el candidato sobrevivió la prueba de esferas bajo los radios, muestras y parámetros usados.
\end{advertencia}

% ============================================================
\section{Workflow 3: \texttt{refined\_basin}}

\subsection{Comando}

\begin{lstlisting}[style=terminal]
hidden-attractors-refined-basin
\end{lstlisting}

O:

\begin{lstlisting}[style=terminal]
python tools/cli/refine_project_basin_classification.py
\end{lstlisting}

Para revisar opciones:

\begin{lstlisting}[style=terminal]
hidden-attractors-refined-basin --help
\end{lstlisting}

\subsection{Propósito}

Este workflow no construye una cuenca desde cero. Revisa una cuenca gruesa previa y toma las celdas marcadas como:

\begin{center}
\texttt{unknown}
\end{center}

Luego reintegra esas condiciones iniciales y compara sus trayectorias contra referencias positiva y negativa mediante:

\begin{itemize}
  \item distancia entre nubes de puntos de cola;
  \item rangos por coordenada;
  \item frecuencia dominante FFT;
  \item sección de Poincaré cuando hay suficientes cruces.
\end{itemize}

\subsection{Cuándo se usa}

Úsalo solamente si ya existe una corrida previa de cuenca con archivos como:

\begin{lstlisting}[style=terminal]
basin_comparison_config.json
project_best_basin_grid.csv
project_best_basin_xy.png
\end{lstlisting}

Si esos archivos no existen, el workflow no tiene fuente que refinar.

\subsection{Ejecución recomendada}

\begin{lstlisting}[style=terminal]
hidden-attractors-refined-basin ^
  --source-dir outputs/basin_compare_danca_project_h001_grid101_20260517 ^
  --output-dir outputs/manual_run/refined_basin_test ^
  --chunks 8
\end{lstlisting}

En Linux/macOS:

\begin{lstlisting}[style=terminal]
hidden-attractors-refined-basin \
  --source-dir outputs/basin_compare_danca_project_h001_grid101_20260517 \
  --output-dir outputs/manual_run/refined_basin_test \
  --chunks 8
\end{lstlisting}

\subsection{Archivos que genera}

\begin{center}
\begin{tabular}{L{0.38\linewidth}L{0.52\linewidth}}
\toprule
Archivo & Significado \\
\midrule
\texttt{refined\_basin\_config.json} & Contrato completo de refinamiento. \\
\texttt{unknown\_refinement\_plan.csv} & Celdas \texttt{unknown} que serán reintegradas. \\
\path{unknown_refined_chunk_*.csv} & Resultados por proceso. \\
\path{unknown_refined_rows.csv} & Todas las celdas refinadas. \\
\path{project_best_basin_refined_grid.csv} & Cuenca final con celdas refinadas. \\
\path{project_best_basin_refined_grid.npy} & Matriz de clases para uso en Python. \\
\path{project_best_basin_refined_xy.png} & Imagen de la cuenca refinada. \\
\path{refined_basin_summary.json} & Resumen final. \\
\bottomrule
\end{tabular}
\end{center}

\subsection{Clases refinadas}

El workflow conserva las clases base y agrega:

\begin{center}
\begin{tabular}{cl}
\toprule
ID & Clase \\
\midrule
0 & \texttt{equilibrium} \\
1 & \texttt{target\_positive} \\
2 & \texttt{target\_negative} \\
3 & \texttt{infinity} \\
4 & \texttt{unknown} \\
5 & \texttt{numerical\_failure} \\
6 & \texttt{bounded\_other} \\
\bottomrule
\end{tabular}
\end{center}

La clase \texttt{bounded\_other} significa que la trayectoria fue acotada y no colapsada, pero no se pareció lo suficiente a las referencias objetivo bajo los umbrales definidos.

% ============================================================
\section{Clasificación de cuencas: significado de las etiquetas}

La librería usa las siguientes etiquetas operativas:

\begin{center}
\begin{tabular}{L{0.12\linewidth}L{0.28\linewidth}L{0.50\linewidth}}
\toprule
ID & Etiqueta & Significado práctico \\
\midrule
0 & \texttt{equilibrium} & La trayectoria termina cerca de un equilibrio. \\
1 & \texttt{target\_positive} & La trayectoria cae en la clase objetivo positiva. \\
2 & \texttt{target\_negative} & La trayectoria cae en la clase objetivo negativa. \\
3 & \texttt{infinity} & La trayectoria diverge o excede umbrales de norma. \\
4 & \texttt{unknown} & El clasificador no decide con los criterios actuales. \\
5 & \texttt{numerical\_failure} & Falló la integración o el cálculo. \\
6 & \texttt{bounded\_other} & Clase agregada en refinamiento: acotada, no colapsada, pero distinta de referencia objetivo. \\
\bottomrule
\end{tabular}
\end{center}

Las clases objetivo son:

\begin{equation}
\{1,2\}=\{\texttt{target\_positive},\texttt{target\_negative}\}.
\end{equation}

\begin{advertencia}
Una etiqueta \texttt{target\_positive} o \texttt{target\_negative} no significa por sí sola ``atractor oculto''. Solo indica que, bajo el contrato numérico probado, la trayectoria fue asignada a una clase objetivo acotada no trivial.
\end{advertencia}

% ============================================================
\section{Cómo leer una corrida completa}

Supón que ejecutaste:

\begin{lstlisting}[style=terminal]
hidden-attractors-sphere-controls --output-dir outputs/manual_run/sphere_controls_deep
\end{lstlisting}

Revisa en este orden:

\subsection{Paso 1: verificar que la corrida terminó}

Abre:

\begin{lstlisting}[style=terminal]
outputs/manual_run/sphere_controls_deep/top3_sphere_robustness_summary.json
\end{lstlisting}

Busca:

\begin{lstlisting}[style=terminal]
"status": "ok"
\end{lstlisting}

Si dice \texttt{partial}, revisa \texttt{logs/}.

\subsection{Paso 2: revisar decisión principal}

Abre:

\begin{lstlisting}[style=terminal]
outputs/manual_run/sphere_controls_deep/sphere_decision.csv
\end{lstlisting}

La columna central es:

\begin{lstlisting}[style=terminal]
hiddenness_status
\end{lstlisting}

\subsection{Paso 3: revisar cuántos contactos hubo}

Si:

\begin{equation}
\texttt{total\_target\_hits}>0,
\end{equation}

entonces hubo contacto entre una vecindad de equilibrio y la clase objetivo. Bajo esa prueba, no debes afirmar ocultedad.

Si:

\begin{equation}
\texttt{total\_target\_hits}=0,
\end{equation}

entonces el candidato sobrevivió los radios y muestras probadas.

\subsection{Paso 4: revisar robustez del propio atractor candidato}

Abre:

\begin{lstlisting}[style=terminal]
outputs/manual_run/sphere_controls_deep/attractor_robustness_summary.csv
\end{lstlisting}

Busca:

\begin{lstlisting}[style=terminal]
robust_attractor
robustness_status
\end{lstlisting}

Si el candidato no es robusto, aunque no haya contactos desde equilibrios, todavía no es un candidato fuerte.

\subsection{Paso 5: revisar crudos si hay dudas}

Abre:

\begin{lstlisting}[style=terminal]
outputs/manual_run/sphere_controls_deep/sphere_raw.csv
\end{lstlisting}

Filtra por:

\begin{lstlisting}[style=terminal]
target_hit = True
\end{lstlisting}

Ahí puedes ver:

\begin{itemize}
  \item equilibrio de origen;
  \item radio;
  \item muestra;
  \item condición inicial exacta;
  \item clase asignada.
\end{itemize}

Esto sirve para reproducir casos críticos.

% ============================================================
\section{Paralelización: qué sí y qué no}

\subsection{Qué sí se paraleliza}

Los workflows paralelizan tareas independientes:

\begin{itemize}
  \item diferentes candidatos en \texttt{robustness\_overlay};
  \item diferentes fragmentos de muestras esféricas en \texttt{sphere\_controls};
  \item diferentes celdas \texttt{unknown} en \texttt{refined\_basin}.
\end{itemize}

\subsection{Qué no debe paralelizarse ingenuamente}

La continuación fraccionaria con memoria no debe romperse como si fuera una colección de condiciones iniciales independientes cuando depende de la historia previa. En sistemas de Caputo, el estado efectivo depende de una ventana de memoria:

\begin{equation}
\{X(t_k-M),\ldots,X(t_k)\}.
\end{equation}

Por eso, si en el futuro se integra continuación en \texttt{version\_2}, debe conservarse la ventana de historia entre pasos de continuación.

\subsection{Uso práctico de \texttt{--chunks}}

En \texttt{sphere\_controls}:

\begin{lstlisting}[style=terminal]
hidden-attractors-sphere-controls --chunks 8
\end{lstlisting}

significa que el plan de condiciones iniciales se divide en 8 subconjuntos independientes. Cada proceso escribe su propio archivo:

\begin{lstlisting}[style=terminal]
sphere_raw_chunk_000.csv
sphere_raw_chunk_001.csv
...
\end{lstlisting}

Después el agregador produce:

\begin{lstlisting}[style=terminal]
sphere_raw.csv
sphere_decision.csv
top3_sphere_robustness_summary.json
\end{lstlisting}

% ============================================================
\section{Uso desde Python como librería}

Además de usar comandos de terminal, puedes escribir tus propios scripts.

\subsection{Ejemplo 1: imprimir equilibrios}

\begin{lstlisting}[style=python]
import numpy as np
from hidden_attractors import chua_piecewise_parameters
from hidden_attractors.models import equilibria_piecewise, rhs_piecewise

params = chua_piecewise_parameters()
for name, eq in equilibria_piecewise(params).items():
    print(name, eq, np.linalg.norm(rhs_piecewise(eq, params)))
\end{lstlisting}

\subsection{Ejemplo 2: cargar candidatos}

\begin{lstlisting}[style=python]
from hidden_attractors import load_final_candidate_records

records = load_final_candidate_records()
for r in records:
    print(r.candidate_id)
    print("route =", r.route)
    print("q =", r.q)
    print("seed =", r.seed)
    print("robust_start =", r.robust_start)
\end{lstlisting}

\subsection{Ejemplo 3: clasificar etiquetas de cuenca}

\begin{lstlisting}[style=python]
from hidden_attractors.basins import class_label, is_target_class

for class_id in range(6):
    print(class_id, class_label(class_id), is_target_class(class_id))
\end{lstlisting}

\subsection{Ejemplo 4: calcular métricas de trayectoria}

\begin{lstlisting}[style=python]
from hidden_attractors.io import load_trajectory_csv
from hidden_attractors.analysis import trajectory_metrics

traj = load_trajectory_csv("outputs/mi_corrida/trayectoria.csv")
metrics, payload = trajectory_metrics(
    traj,
    h=0.01,
    t_start=750.0,
)

for k, v in metrics.items():
    print(k, v)
\end{lstlisting}

% ============================================================
\section{Relación con función descriptiva y Machado/FDF}

La librería \texttt{version\_2} carga candidatos que provienen de rutas anteriores, incluyendo ruta Lur'e clásica y ruta Machado/FDF. Operativamente:

\begin{itemize}
  \item la función descriptiva clásica y la función descriptiva fraccionaria se usan como generadores heurísticos de semillas;
  \item esas semillas no prueban ciclos exactos en Caputo;
  \item las semillas deben trasladarse al sistema no lineal completo mediante simulación y verificación;
  \item \texttt{version\_2} se concentra en verificar, comparar, clasificar y documentar esos candidatos.
\end{itemize}

\subsection{Puente Weyl--Caputo}

Cuando se usan Fourier, balance armónico, Nyquist o función descriptiva, el razonamiento armónico ideal corresponde al régimen estacionario tipo Weyl/Liouville--Weyl. La simulación causal de la librería corresponde a Caputo, con memoria desde $t_0$. La diferencia aparece como discrepancia de memoria. Por tanto, las semillas armónicas se interpretan como aproximaciones útiles para explorar el sistema de Caputo, no como prueba de ciclos periódicos exactos.

% ============================================================
\section{Errores comunes}

\subsection{Error: \texttt{ModuleNotFoundError: hidden\_attractors}}

Causa probable: no instalaste el paquete o no estás en el entorno virtual correcto.

Solución:

\begin{lstlisting}[style=terminal]
cd version_2
python -m pip install -e .
\end{lstlisting}

\subsection{Error: no encuentra archivos en \texttt{outputs/}}

Algunos cargadores esperan salidas de referencia. Verifica que exista la carpeta esperada. Por ejemplo, los candidatos finales pueden depender de archivos en:

\begin{lstlisting}[style=terminal]
outputs/extended_search/
outputs/lure_biased_multiparam_q09998_20260515_195444/
\end{lstlisting}

Si no existen, el cargador no puede inventar candidatos.

\subsection{Error de compilador C}

Si falla un workflow nativo, verifica que tengas compilador C.

En PowerShell:

\begin{lstlisting}[style=terminal]
gcc --version
\end{lstlisting}

Si no existe, instala MinGW, MSYS2 o un compilador compatible y agrégalo al \texttt{PATH}.

\subsection{La corrida parece no terminar}

Revisa los logs:

\begin{lstlisting}[style=terminal]
outputs/tu_corrida/logs/
\end{lstlisting}

Los archivos \texttt{.out} muestran avance. Los archivos \texttt{.err} muestran errores.

\subsection{El resumen dice \texttt{partial}}

Significa que no todos los procesos terminaron o no todos escribieron sus archivos \texttt{.done}. Revisa:

\begin{itemize}
  \item archivos \texttt{.done};
  \item logs de errores;
  \item si se interrumpió la terminal;
  \item si hubo error de compilación nativa;
  \item si la ruta \texttt{--source-dir} existe.
\end{itemize}

% ============================================================
\section{Checklist para defender una corrida}

Antes de usar resultados en un reporte, asegúrate de poder llenar esta lista:

\begin{enumerate}[label=\arabic*]
  \item ¿Qué modelo se usó? Escribir ecuaciones y parámetros.
  \item ¿Qué orden fraccionario $q$ se usó?
  \item ¿Qué método de integración se usó? EFORK nativo, ABM u otro.
  \item ¿Qué paso $h$ se usó?
  \item ¿Qué memoria finita $L_m$ se usó?
  \item ¿Cuál fue el tiempo final?
  \item ¿Cuál fue el tiempo transitorio descartado?
  \item ¿Cuáles son todos los equilibrios?
  \item ¿Desde qué candidato o semilla se inició la trayectoria objetivo?
  \item ¿La trayectoria candidata fue robusta al cambiar $h$, $L_m$ y $T$?
  \item ¿Se lanzaron trayectorias desde vecindades de todos los equilibrios?
  \item ¿Qué radios se probaron?
  \item ¿Cuántas muestras por radio se usaron?
  \item ¿Hubo \texttt{target\_hits} desde algún equilibrio?
  \item ¿Qué archivo CSV/JSON respalda cada afirmación?
\end{enumerate}

\subsection{Frases correctas para el reporte}

Si no hubo impactos desde equilibrios:

\begin{quote}
Bajo el contrato numérico especificado, las pruebas esféricas alrededor de los equilibrios no detectaron intersección entre las vecindades probadas y la cuenca objetivo. Por tanto, el candidato es compatible con ocultedad bajo los radios, muestras y parámetros evaluados.
\end{quote}

Si sí hubo impactos:

\begin{quote}
La prueba esférica detectó trayectorias iniciadas en vecindades de equilibrio que alcanzan la clase objetivo. Bajo este contrato numérico, la hipótesis de atractor oculto no queda respaldada para este candidato.
\end{quote}

Si solo se hizo \texttt{robustness\_overlay}:

\begin{quote}
La trayectoria candidata muestra persistencia geométrica y espectral bajo variaciones de $h$, $L_m$ y tiempo final. Este resultado evalúa robustez numérica, pero no constituye una verificación de ocultedad.
\end{quote}

% ============================================================
\section{Ruta de trabajo sugerida para tus próximas corridas}

Para una nueva sesión de trabajo, usa esta secuencia:

\subsection*{Día 1: verificación rápida}

\begin{lstlisting}[style=terminal]
cd C:\Users\moren\Desktop\Hidden-Attractors-Localization\version_2
.\.venv\Scripts\Activate.ps1
python -m compileall hidden_attractors examples tests tools/cli
python examples/quickstart_equilibria.py
python examples/list_final_candidates.py
\end{lstlisting}

\subsection*{Día 1: robustez ligera}

\begin{lstlisting}[style=terminal]
hidden-attractors-robustness-overlay --output-dir outputs/manual_run/robustness_overlay_light
\end{lstlisting}

Revisar:

\begin{lstlisting}[style=terminal]
outputs/manual_run/robustness_overlay_light/robustness_overlay_summary.json
outputs/manual_run/robustness_overlay_light/robustness_overlay_metrics.csv
outputs/manual_run/robustness_overlay_light/plots/
\end{lstlisting}

\subsection*{Día 2: prueba de ocultedad por esferas}

\begin{lstlisting}[style=terminal]
hidden-attractors-sphere-controls ^
  --output-dir outputs/manual_run/sphere_controls_main ^
  --samples-per-radius 500 ^
  --chunks 6
\end{lstlisting}

Revisar:

\begin{lstlisting}[style=terminal]
outputs/manual_run/sphere_controls_main/sphere_decision.csv
outputs/manual_run/sphere_controls_main/sphere_cumulative_summary.csv
outputs/manual_run/sphere_controls_main/attractor_robustness_summary.csv
\end{lstlisting}

\subsection*{Día 3: refinamiento si hay cuenca gruesa}

Solo si existe una cuenca previa:

\begin{lstlisting}[style=terminal]
hidden-attractors-refined-basin ^
  --source-dir outputs/basin_compare_danca_project_h001_grid101_20260517 ^
  --output-dir outputs/manual_run/refined_basin_main ^
  --chunks 8
\end{lstlisting}

% ============================================================
\section{Mapa de módulos principales}

\begin{center}
\begin{tabular}{L{0.38\linewidth}L{0.52\linewidth}}
\toprule
Módulo & Para qué sirve \\
\midrule
\path{hidden_attractors.models.chua} & Define parámetros, no linealidad, campo vectorial y equilibrios del Chua no suave. \\
\path{hidden_attractors.candidates} & Carga candidatos finales desde salidas de referencia. \\
\path{hidden_attractors.analysis.trajectory} & Calcula métricas de trayectoria: rangos, varianzas, FFT, secciones, distancias entre nubes. \\
\path{hidden_attractors.basins.classification} & Define etiquetas de clases de cuenca. \\
\path{hidden_attractors.io} & Lee y escribe CSV/JSON, carga trayectorias y serializa resultados. \\
\path{hidden_attractors.plotting} & Genera figuras de fase, proyecciones, series, bifurcaciones y overlays. \\
\path{hidden_attractors.native} & Envoltorios hacia integración EFORK y clasificación nativa en C. \\
\path{hidden_attractors.workflows.robustness_overlay} & Workflow de robustez geométrica y espectral. \\
\path{hidden_attractors.workflows.sphere_controls} & Workflow de controles alrededor de equilibrios para evaluar ocultedad. \\
\path{hidden_attractors.workflows.refined_basin} & Workflow de refinamiento de celdas \texttt{unknown} en cuencas. \\
\bottomrule
\end{tabular}
\end{center}

% ============================================================
\section{Resumen final}

La librería se usa en tres niveles:

\begin{enumerate}[label=\arabic*]
  \item \textbf{Nivel básico:} instalar, verificar equilibrios, listar candidatos y cargar trayectorias.
  \item \textbf{Nivel operativo:} ejecutar workflows de robustez, controles esféricos y refinamiento de cuencas.
  \item \textbf{Nivel de investigación:} interpretar los resultados bajo el criterio de atractores ocultos, cuidando que cada afirmación esté respaldada por el contrato numérico y sus archivos de salida.
\end{enumerate}

La ruta mínima para una verificación defendible es:

\begin{equation}
\text{candidato} \longrightarrow \text{robustez} \longrightarrow \text{controles alrededor de equilibrios} \longrightarrow \text{lectura de cuencas}.
\end{equation}

No debe concluirse ocultedad solo desde una semilla de función descriptiva, ni solo desde una trayectoria atractora, ni solo desde robustez visual. La prueba numérica relevante es la ausencia de contacto detectado entre la cuenca objetivo y vecindades de todos los equilibrios bajo radios, muestras y parámetros explícitos.

% ============================================================
\appendix
\section{Comandos de referencia rápida}

\subsection*{Instalar}

\begin{lstlisting}[style=terminal]
cd version_2
python -m pip install -e ".[dev]"
\end{lstlisting}

\subsection*{Probar}

\begin{lstlisting}[style=terminal]
python -m compileall hidden_attractors examples tests tools/cli
python examples/quickstart_equilibria.py
python examples/list_final_candidates.py
python -m pytest -q
\end{lstlisting}

\subsection*{Ayuda de workflows}

\begin{lstlisting}[style=terminal]
hidden-attractors-robustness-overlay --help
hidden-attractors-sphere-controls --help
hidden-attractors-refined-basin --help
\end{lstlisting}

\subsection*{Robustez}

\begin{lstlisting}[style=terminal]
hidden-attractors-robustness-overlay --output-dir outputs/manual_run/robustness_overlay
\end{lstlisting}

\subsection*{Controles esféricos}

\begin{lstlisting}[style=terminal]
hidden-attractors-sphere-controls --output-dir outputs/manual_run/sphere_controls --chunks 8 --samples-per-radius 1000
\end{lstlisting}

\subsection*{Refinamiento de cuenca}

\begin{lstlisting}[style=terminal]
hidden-attractors-refined-basin --source-dir outputs/basin_compare_danca_project_h001_grid101_20260517 --output-dir outputs/manual_run/refined_basin --chunks 8
\end{lstlisting}

\section{Archivos que conviene guardar para el reporte}

\begin{itemize}
  \item \texttt{robustness\_overlay\_config.json}
  \item \texttt{robustness\_overlay\_metrics.csv}
  \item \texttt{robustness\_overlay\_summary.json}
  \item \texttt{plots/overlay\_*.png}
  \item \texttt{top3\_sphere\_robustness\_config.json}
  \item \texttt{sphere\_decision.csv}
  \item \texttt{sphere\_cumulative\_summary.csv}
  \item \texttt{sphere\_raw.csv}
  \item \texttt{attractor\_robustness\_summary.csv}
  \item \texttt{refined\_basin\_summary.json}, si se ejecutó refinamiento
  \item \texttt{project\_best\_basin\_refined\_xy.png}, si se ejecutó refinamiento
\end{itemize}

\end{document}
